\documentclass[11pt]{letter}

\usepackage{geometry}
\geometry{margin=1in}
\usepackage{setspace}
\usepackage{hyperref}

\signature{Decory Edwards}
\address{Decory Edwards \\ Johns Hopkins University \\ Department of Economics}

\begin{document}

\begin{letter}{Hiring Committee \\ The Voleon Group \\ Berkeley, CA}

\opening{Dear Hiring Committee,}

I am writing to apply for the Member of Research Staff position at Voleon. I am currently a PhD candidate in Economics at Johns Hopkins University. My training combines mathematical reasoning, statistical modeling, and data driven empirical analysis. I am motivated by problems that require precise logic, careful implementation, and validation against noisy real world data. Voleon’s focus on applying modern machine learning methods to financial markets aligns closely with how I approach research.

My research agenda is at the intersection of wealth inequality and macroeconomics. In my job market paper, I build and estimate a structural heterogeneous agent model with returns that differ across households. The core contribution is to show how heterogeneity in returns and income risk jointly shape the wealth distribution and consumption behavior. Relative to closely related work, my framework performs better at matching the lower and middle portions of the wealth distribution. This difference matters quantitatively. Models that focus primarily on returns heterogeneity can fit the top of the wealth distribution closely but tend to generate too much wealth among the bottom half of households. In my framework, income uncertainty generates a precautionary saving motive that produces genuinely low wealth households. This leads to a realistic distribution of marginal propensities to consume and aggregate consumption responses that align with empirical estimates.

A key feature of my research is the heavy use of computation to solve and simulate models that cannot be handled analytically. I write code to solve dynamic programming problems, simulate large populations, and evaluate model performance under uncertainty. I work with complex and diverse datasets and design experiments to assess predictive content and robustness. This work requires careful numerical reasoning, attention to statistical assumptions, and disciplined validation, skills that are directly applicable to developing and testing predictive models in financial markets.

I am particularly interested in Voleon because of its emphasis on applied research across the full life cycle, from idea generation to production deployment. The opportunity to work alongside researchers and engineers to translate research insights into live trading systems is especially appealing. I value environments that combine academic rigor with immediate feedback and where research impact is measured through real world performance.

My economics training has provided a strong foundation in probability, optimization, and statistical inference. I am comfortable reasoning from first principles, translating theory into working models, and writing research grade and production oriented code in Python. I also value collaboration and clear communication, both of which are essential when research outcomes must integrate with large scale systems.

I am also enthusiastic about the prospect of living and working in Berkeley and participating in a hybrid research environment. I value in person collaboration while also appreciating focused research time. I see this role as an opportunity to apply my academic training in a setting where results are immediate and meaningful.

I would welcome the opportunity to contribute my skills and perspective to Voleon. Thank you for your time and consideration.

\closing{Sincerely,}

\end{letter}
\end{document}
